\subsection{A binary-degree-twelve parameter construction correcting a divisibility claim}
\label{sec:binary12-hsop}

\paragraph{P-F43 (a parameter system with no degree divisible by ten).}
Let
\[
 R=\mathbb C[\operatorname{Sym}^{12}\mathbb C^2]^{SL_2},
 \qquad X=\operatorname{Spec}R,
 \qquad A=5544=2^3\,3^2\,7\,11,
 \qquad B=3465=3^2\,5\,7\,11.
\]
For every sufficiently large integer $k\equiv1\pmod {10}$,
there is a homogeneous system of parameters of $R$ with degrees
\[
 \underbrace{Ak,Ak,Ak,Ak,Ak}_{5\text{ entries}},\quad
 \underbrace{Bk,Bk,Bk,Bk,Bk}_{5\text{ entries}}.
 \tag{P-F43a}
\]
One explicit valid choice is $k=10^{12}+1$, giving five copies
of each of the degrees
\[
 a_0=5\,544\,000\,000\,005\,544,\qquad
 b_0=3\,465\,000\,000\,003\,465.
 \tag{P-F43a'}
\]
None of these ten degrees is divisible by ten. This contradicts
the literal assertion of Brouwer--Draisma--Popoviciu,
\emph{The degrees of a system of parameters of the ring of
invariants of a binary form}, Theorem~3.4(ii), printed p.~955,
at binary order $12$ and $t=10$. It also contradicts the
degree-ten necessity stated in the Note on Brouwer's
Poincar\'e-series page. The result does not decide whether
that page's much smaller tuple
$(2,4,5,5,6,6,7,8,9,11)$ is the degree sequence of a parameter
system. The construction supplies explicit valid degrees,
but not explicit polynomial parameters or an efficient
procedure for constructing them.

\paragraph{Proof: uniform eventual nonvanishing from finite generation.}
The grading defines an action of $\mathbb C^*$ on $X$, whose
vertex is denoted $0$. For $x\ne0$, set
\[
 S_x=\{d\ge0:\text{some }f\in R_d\text{ satisfies }f(x)\ne0\},
 \qquad g_x=\gcd(S_x\setminus\{0\}).
 \tag{P-F43b}
\]
Choose finitely many homogeneous positive-degree algebra
generators $h_1,\ldots,h_N$ of $R$. There is at least one
$h_i(x)\ne0$ for every $x\ne0$. The semigroup $S_x$ is exactly
the additive semigroup generated by the degrees of the
$h_i$ which are nonzero at $x$. One inclusion follows by
multiplying these nonzero values. For the other, express a
homogeneous $f$ as a weighted-homogeneous polynomial in the
$h_i$; a nonzero evaluation has at least one nonzero monomial
term, and that term has the same degree as $f$.

Only finitely many semigroups $S_x$ therefore occur. Each
$S_x/g_x$ is a finitely generated subsemigroup of
$\mathbb Z_{\ge0}$ with gcd one and contains all sufficiently
large integers. For completeness, choose a positive element
$a$ of such a semigroup. Its residues modulo $a$ form a
subsemigroup of the finite group $\mathbb Z/a$, hence a
subgroup, and gcd one makes it the whole group. Choose one
semigroup representative of each residue; adding multiples
of $a$ proves eventual membership. Taking a maximum over
the finitely many semigroups gives an integer $C$ such that
\[
 d\ge C\quad\Longrightarrow\quad
 \bigl(d\in S_x\ \Longleftrightarrow\ g_x\mid d\bigr)
 \qquad\text{for every }x\in X\setminus\{0\}.
 \tag{P-F43c}
\]
Moreover the stabilizer of $x$ for the grading action is
exactly $\mu_{g_x}$: a scalar $\lambda$ fixes $x$ if and only
if $\lambda^d=1$ for every degree with a nonzero value there.
In particular every nonvertex grading orbit has dimension one.

\paragraph{Scalar orders on the binary quotient.}
The quotient map
\[
 \pi:\operatorname{Sym}^{12}\mathbb C^2\longrightarrow X
\]
has one closed $SL_2$-orbit in each fiber. The nonzero closed
orbits are represented either by a stable binary form, whose
root multiplicities are all less than six, or by
$c x^6y^6$ with $c\ne0$. These are the standard binary-form
stability and closed-orbit criteria; see Alper--Isaev,
\emph{Associated forms and hypersurface singularities: the
binary case}, Section~5.2, printed p.~13, paragraph preceding
Proposition~5.1. In the strictly semistable case the affine
description also follows directly: put a multiplicity-six
root at $x=0$ and apply the diagonal one-parameter subgroup
which keeps the nonzero $x^6y^6$ term and sends all the other
terms to zero. A closed orbit must contain this limit.

Suppose first that a stable form $F$ represents $x$. It has
at least three distinct roots, so its projective automorphism
group is finite: it permutes that finite set, and a projective
linear transformation fixing three distinct points is the
identity. Choose a primitive $g_x$th root $\lambda$.
The closed orbits of $\lambda F$ and $F$ have the same
quotient point, so some $S\in SL_2$ satisfies
$S\cdot F=\lambda F$. Its projective class has finite order.
Consequently $S$ is diagonalizable, and after an $SL_2$
change of coordinates it has the form
$\operatorname{diag}(z,z^{-1})$. Write
\[
 t=\operatorname{ord}(z^2),\qquad
 F=\sum_{i=0}^{12}c_i x^{12-i}y^i.
\]
All nonzero coefficients have indices in a single residue
class $j$ modulo $t$. Since $F$ is not a monomial, it has
two distinct such indices, and therefore $t\le12$.
In the convention where the indicated diagonal matrix acts
on the displayed monomials with weight $z^{12-2i}$, one has
$\lambda=(z^2)^{6-j}$; the inverse action convention gives
the same order. Thus
\[
               g_x=\frac{t}{\gcd(t,6-j)},\qquad 1\le t\le12.
 \tag{P-F43d}
\]
If $10\mid g_x$, then $g_x=t=10$. Two supported indices
modulo ten force $j\in\{0,1,2\}$. Formula (P-F43d) gives
orders $5,2,5$, respectively, a contradiction.

For $F=cx^6y^6$, every projective transformation taking $F$
to a scalar multiple preserves its unordered pair of roots.
An $SL_2$ lift is diagonal or antidiagonal. A diagonal lift
fixes $F$, and an antidiagonal lift multiplies it by
$(-1)^6=1$. Hence $g_x=1$ in this case. We have proved
\[
              1\le g_x\le12,\qquad 10\nmid g_x
                    \quad(x\in X\setminus\{0\}).
 \tag{P-F43e}
\]

\paragraph{Dimensions of the two exceptional loci.}
Let
\[
 E_2=X^{\mu_2}=\{0\}\cup\{x\ne0:2\mid g_x\},
 \qquad
 E_5=X^{\mu_5}=\{0\}\cup\{x\ne0:5\mid g_x\}.
\]
These are closed conical subsets, and (P-F43e) gives
$E_2\cap E_5=\{0\}$. We claim
\[
                       \dim E_2\le5,\qquad\dim E_5\le2.
 \tag{P-F43f}
\]
For $1\le t\le12$ and $0\le j<t$, let $U_{t,j}$ be the
linear space of binary forms whose supported indices are
congruent to $j$ modulo $t$. The preceding diagonalization
places every point in question in the image of the stable
part of one of finitely many such spaces. The diagonal
$SL_2$ torus preserves $U_{t,j}$ and has one-dimensional
orbits on its stable part, because a stable form has finite
projective stabilizer. Consequently
\[
 \dim\pi(U_{t,j}\cap\mathrm{stable})
       \le \dim U_{t,j}-1.
 \tag{P-F43g}
\]
If $g_x$ is even, $t$ is even. For $t=2$, (P-F43d) forces
$j=1$; there are six odd indices between zero and twelve,
so (P-F43g) bounds the image dimension by five. For $t\ge3$
there are at most $\lfloor12/t\rfloor+1\le5$ indices in any
residue class, giving image dimension at most four. If
$5\mid g_x$, then $t=5$ or $10$ and each residue class has
at most three indices, so its image dimension is at most
two. Strictly semistable closed representatives contribute
to neither locus. Taking the finite union and its closure
proves (P-F43f).

\paragraph{The entire degree spaces have complementary base loci.}
Choose $k\equiv1\pmod {10}$ sufficiently large that
$a=Ak$ and $b=Bk$ are at least $C$. For each integer
$g\le12$, one has
\[
 g\mid a\ \Longleftrightarrow\ 5\nmid g,
 \qquad
 g\mid b\ \Longleftrightarrow\ 2\nmid g.
\]
Indeed $A$ contains every prime power at most twelve except
five, and $k$ is not divisible by five; $B$ contains every
odd prime power at most twelve, and $k$ is odd. By
(P-F43c), the common zero locus in $X$ of the entire space
$R_a$ is $E_5$, and that of $R_b$ is $E_2$.

Since $\dim X=12-2=10$, a nonempty subsequence of at most
five $a$-degrees has base locus of codimension at least
eight, a subsequence of at most five $b$-degrees has base
locus of codimension at least five, and any mixed
subsequence has base locus $\{0\}$ of codimension ten.
These are precisely the hypotheses of Dixmier's criterion,
as stated in Brouwer--Draisma--Popoviciu, Proposition~5.1,
printed p.~958 (preprint p.~6). They already imply (P-F43a).

Here is an independent incidence proof of this last
selection step. Let $\mathcal V=R_a^5\times R_b^5$ be the
space of parameter choices. Stratify $X\setminus\{0\}$
into $E_5\setminus\{0\}$, $E_2\setminus\{0\}$, and their
complement. For a fixed point of the first stratum, the
condition that all ten choices vanish imposes five
independent linear conditions on $\mathcal V$; the other
five vanish automatically. The associated incidence
variety has dimension at most $\dim\mathcal V+2-5$.
For the second stratum the corresponding bound is
$\dim\mathcal V+5-5$, and for the complement it is
$\dim\mathcal V+10-10$. These statements follow equally
by viewing the incidence as the kernel of the evaluation
map of constant rank over each stratum.

Each of these incidence varieties is invariant under
the grading $\mathbb C^*$ action on the point coordinate,
with the parameter choice held fixed: all ten equations
are homogeneous. Its irreducible components are preserved
by this connected group. Every nonempty fiber of an
irreducible component over $\mathcal V$ contains a
one-dimensional grading orbit. The dimension of its image
is therefore at most the dimension of the component minus
one. All the bad images have dimension strictly less than
$\dim\mathcal V$. There are only finitely many components,
so their closures cannot cover $\mathcal V$. A generic
choice has the vertex as its entire common zero set.

For such a choice $f_1,\ldots,f_{10}$, the graded algebra
$R/(f_1,\ldots,f_{10})$ is finite-dimensional by the
Nullstellensatz. Lifting a finite homogeneous vector-space
basis and inducting on degree shows that $R$ is finite as
a module over $\mathbb C[f_1,\ldots,f_{10}]$. This subring
has dimension ten because a finite integral extension
preserves dimension. It is generated by ten elements and
is a domain, so these elements are algebraically
independent: a nonzero prime kernel in the polynomial
ring on ten variables would lower its dimension. They
are therefore a homogeneous system of parameters.
Finally $5\nmid a$ and $2\nmid b$, proving the assertion
about divisibility by ten.\hfill$\square$

\paragraph{An explicit uniform threshold and degree tuple.}
To verify the advertised value $k=10^{12}+1$, use Derksen,
\emph{Polynomial bounds for rings of invariants},
Theorem~1.1, Proposition~1.2, and Example~2.1, printed
p.~956. In characteristic zero these give for binary
forms of degree $d$ an algebra-generator degree bound
$\frac32(d+1)d^6$. The representation of $SL_2$ has finite
kernel, as required by the cited proposition, and its
matrix coefficients have degree at most $d$ in the
matrix entries. At $d=12$ we may therefore choose the
generators above with degrees at most
\[
              \mathcal B=\frac32\cdot13\cdot12^6
                        =58\,226\,688.
 \tag{P-F43h}
\]

Here is the needed elementary uniform conductor estimate.
If an additive semigroup is generated by positive integers
at most $\mathcal B$ and has gcd $g$, divide its generators
by $g$ and let $h$ be the smallest normalized generator.
In the directed graph of residues modulo $h$, with an
edge for addition of each normalized generator, every
residue is reachable from zero. A shortest such path
repeats no vertex, so has at most $h-1$ edges. Hence each
residue has a semigroup representative at most
\[
            (h-1)\mathcal B/g\le(\mathcal B/g)^2.
\]
Every larger integer in the same residue belongs to the
normalized semigroup, by adding copies of $h$.
If $N\ge\mathcal B^2$ and $g\mid N$, then
$N/g\ge\mathcal B^2/g\ge(\mathcal B/g)^2$.
Thus the original semigroup contains every $g$-multiple
at least $\mathcal B^2$. This proof applies uniformly to
all the subsets of generators nonzero at a point, so
$C=\mathcal B^2$ is valid in (P-F43c).

Exact integer arithmetic gives
\[
 \begin{split}
 \mathcal B^2&=3\,390\,347\,195\,449\,344,\\
 3465(10^{12}+1)&=3\,465\,000\,000\,003\,465
                        >\mathcal B^2.
 \end{split}
\]
The other degree is larger still, and $10^{12}+1\equiv1$
modulo ten. The proof just given therefore applies to
the concrete tuple (P-F43a'). The only arithmetic
computation here checks displayed integer values; the
existence assertion follows from the full argument above.

\paragraph{Exact source discrepancy and scope.}
Theorem~3.4(ii) asserts at least $\lfloor(n-1)/t\rfloor$
degrees divisible by $t$ for every integer $t>1$ and every
even binary order $n$. Its preceding Lemma~3.3(ii) gives
instead $\lfloor(n-j)/t\rfloor$, where $j$ is the least
nonnegative integer with $\gcd(n/2-j,t)=1$.
For $(n,t)=(12,10)$ this least $j$ is three, so the lemma
only gives zero. The printed improvement writes $n=at+b$
and replaces a condition on $\gcd(n-2i,2t)$ by one on
$\gcd(b-2i,2t)$. This replacement is invalid when $a$ is
odd: at $n=12,t=10,a=1,b=2,i=0$ the two gcds are four
and two. The construction above proves that this is a
defect in the stated necessity, not merely an unfilled
proof step. It makes no priority claim and does not
invalidate other independently proved results in that paper.

The exact primary sources used here are
\begin{itemize}
\item Brouwer--Draisma--Popoviciu, Theorem~3.4(ii), p.~955,
and Proposition~5.1, p.~958:
\url{https://pure.tue.nl/ws/files/17749340/Degrees2015.pdf};
the latter is Proposition~5.1, p.~6, in
\url{https://arxiv.org/pdf/1404.5722}.
\item Alper--Isaev, Section~5.2, p.~13, paragraph preceding
Proposition~5.1 (binary stability and the strictly semistable
closed orbit):
\url{https://maths-people.anu.edu.au/~alperj/papers/binary.pdf}.
\item Brouwer's Poincar\'e-series page, Note paragraph
following the two displayed expressions for binary order
twelve:
\url{https://aeb.win.tue.nl/math/poincare.html}.
\item Derksen, \emph{Polynomial bounds for rings of
invariants}, Theorem~1.1, Proposition~1.2, and Example~2.1,
p.~956:
\url{https://sites.lsa.umich.edu/hderksen/wp-content/uploads/sites/614/2018/05/A.I.a.7.pdf}.
\end{itemize}
